% =============================================================================
% Abstrak dan Abstract -- keduanya disajikan pada satu halaman.
% =============================================================================

\ipbprelimheading{ABSTRAK}
% Ganti panduan pada abstract.tex dengan narasi abstrak satu paragraf.
\noindent
\MakeUppercase{\NamaPenulis}. \JudulSkripsi\unskip. Dibimbing oleh
\MakeUppercase{\NamaPembimbingSatu} dan \MakeUppercase{\NamaPembimbingDua}


\ipbprelimsubheading{\itshape ABSTRACT}
% Replace the guidance in abstract.tex with a one-paragraph abstract.
\begin{otherlanguage}{english}
\itshape
\noindent
\MakeUppercase{\NamaPenulisInggris}. \JudulSkripsiInggris\unskip. Supervised by
\MakeUppercase{\NamaPembimbingSatu} and \MakeUppercase{\NamaPembimbingDua}
\end{otherlanguage}


\vspace{0.5cm}

Narasi disusun dalam satu paragraf, isi tidak lebih dari 200 kata, dan ditulis
dalam satu halaman untuk abstrak/\textit{abstract}. Abstrak memuat latar
belakang permasalahan (tentatif), tujuan tugas akhir, metode, hasil dengan
penekanan pada temuan baru, dan implikasi yang disajikan secara informatif dan
faktual. Tidak diperbolehkan mengacu pustaka, gambar, dan tabel. Singkatan hanya
dikenalkan jika masih digunakan lagi dalam bagian lain dari
abstrak/\textit{abstract}. Abstrak dalam bahasa Inggris ditulis dengan huruf
miring (\textit{italic}).

\vspace{0.48cm}

\noindent\hangindent=2cm\hangafter=1
Kata kunci: \KataKunci

\begin{otherlanguage}{english}
\itshape
\noindent\hangindent=2cm\hangafter=1
Keywords: \Keywords
\par
\end{otherlanguage}
